% !PRINTED 120
\weberpage{120}
We decompose $P$ as in (4) into the $j$ subsidiary groups $B$, denote
by $e$ all the elements of $Q$, and now investigate the permutation
\[
\pi_e = (B,\ Be),
\]
which is called forth by the elements $e$ among the $j$ elements (4).
By (5) and (6), these permutations certainly form a group $\Pi$
simply or multiply isomorphic with $Q$.

If $a_1$ is any element of $P$, then in the complex
\[
P_1 = Ra_1Q
\]
formed according to the composition of the parts, a certain number of the
subsidiary groups $B$ will be contained. We find this
number by seeking the condition that
$Ra_1 = Ra_1e$, or that $a_1e$ be contained in $Ra_1$, or $e$ in $a_1^{-1}Ra_1$.
If, then,
\[
Q_1 \text{ the intersection of } Q \text{ and } a_1^{-1}Ra_1
\]
is of index $h_1$ with respect to $Q$, then $P_1$ contains exactly $h_1$ and not
more of the complexes $B$.

If $a_2$ is a second element of $P$, then the complex
\[
P_2 = Ra_2Q
\]
either coincides entirely with $P_1$, or the two contain
no common element at all. For if $P_1$ and $P_2$ contain a
common element, then they have one of the complexes $B$ in common. But that is possible only if $a_2e$ is contained in $Ra_1e'$,
and hence $a_2$ in $P_1$; and then $P_1$ and $P_2$ are identical.
We can therefore continue forming the complexes $P_1, P_2 \ldots$ until the group $P$ is exhausted, and obtain the decomposition of the group $P$ with respect to the two groups $Q$ and $R$:
\weq{9}{\begin{aligned}
P &= Ra_1Q + Ra_2Q + \cdots\\
&= P_1 \phantom{{}={}} + P_2 \phantom{{}={}} + \cdots
\end{aligned}}

From the remark that, if $a_2$ does not occur in $Ra_1Q$,
then $a_2^{-1}$ is not contained in $Qa_1^{-1}R$, there results the second
decomposition:
\weq{10}{P = Qa_1^{-1}R + Qa_2^{-1}R + \cdots\footnotemark.}

The parts $P_1, P_2, \ldots$ of (9) contain respectively $h_1, h_2, \ldots$
elements of the sequence $B$, if $h_1, h_2, \ldots$ denote the indices of the intersections of $Q$ with $a_1^{-1}Ra_1, a_2^{-1}Ra_1, \ldots$ with respect to $Q$, and hence % [sic: Weber prints a_2^{-1}Ra_1; the conjugate pattern requires a_2^{-1}Ra_2.]
divisors of the degree $q$ of $Q$.
\footnotetext{I owe these decompositions to a communication from \gesperrt{Dedekind}.}
